\section{Fungsi Proposisional (Kalimat Terbuka)}

Di ranah Logika Proposisi, kepingan huruf seperti $p, q, r$ mewakili secara utuh satu kalimat tunggal deklaratif yang kebetulan memiliki sifat absolut \textbf{sudah tertutup} (memiliki 1 entitas pasti nilai kebenaran). Namun, begitu kita mencampuradukkan Term dari genus Variabel Anonim ke dalam pelukan Predikat, kita melahirkan wujud \textbf{Fungsi Proposisional} \cite{rosen2019}.

\subsection{Definisi Klasik Fungsi Proposisional}

\begin{definisi}[Fungsi Proposisional]
\logic{Fungsi Proposisional} (atau sering diekuivalensikan dalam bahasa linguistik awam sebagai *Kalimat Terbuka*) adalah perpaduan simbiosis mutlak antara sesosok Predikat dengan sekurang-kurangnya sebuah (atau lebih) Variabel tak teridentifikasi bebas ($x, y, z$). 
\end{definisi}

Konsekuensi paling kritis dari definisi ini adalah: \textbf{Sebuah Fungsi Proposisional TIDAK memiliki dan dipastikan gagal dievaluasi nilai kebenarannya (True/False)}. Ini disebabkan karena subjek/aktor yang ingin dievaluasi oleh predikat masih menyembunyikan identitas aslinya di balik topeng variabel (anonim). Nilai kebenarannya masih ``ngambang'' bergantung pada identitas asli dari huruf $x$ kelak.

\subsection{Evolusi dari Kalimat Terbuka Menuju Tertutup}

\begin{contoh}
Misalkan diberikan contoh kalimat verbal: \textit{"Dia adalah Presiden Pertama Republik Indonesia"}. 

Dalam frasa di atas, kata ganti peranti tunjuk \textbf{"Dia"} menduduki peran sentral Variabel bebas ($x$). Di sisi berlawanan, susunan frasa verba klaim \textbf{"...adalah Presiden Pertama Republik Indonesia"} ditahbiskan menduduki fungsi Predikat sentral, dinotasikan kapital $P$. 

Sehingga struktur rangka kerangka aslinya hanyalah: $P(\text{Dia})$ atau lebih lazim ditulis secara matematis akademis menjadi $P(x)$. 

Pada wujud mentah seperti ini, kita jelas takkan bernyali menjatuhkan palu hakim putusan True atau False. Kita tak tahu siapa si \textbf{x} tersebut. Bentuk luntang-lantung inilah yang disahkan sebagai \textbf{Kalimat Terbuka (Fungsi Proposisional)}.
\end{contoh}

\subsection{Mensubstitusi dengan Konstanta Absolut}

Dari statusnya yang masih menjadi pergunjingan ambigu, Fungsi Proposisional $P(x)$ dapat menetas bermetamorfosis menjadi sebuah \textbf{Proposisi Sempurna} (Kalimat Tertutup) hanya dan hanya jika kita mengeksekusi proses **Substitusi Identitas Konstanta** pada variabel pengembara $x$ tersebut.

\begin{contoh}
Melanjutkan contoh $P(x)$ sebelumnya. Mari kita tembak variabel $x$ secara membabi buta menggunakan 3 sampel selongsong uji Konstanta berbeda:

\begin{enumerate}
    \item \textbf{Subjek Konstanta Pertama:} $a = \text{B.J. Habibie}$
    \begin{itemize}
        \item Substitusi ke wajan predikat: $P(a)$
        \item Kalimat merangkai utuh menjadi: \textit{"B.J. Habibie adalah Presiden Pertama Republik Indonesia."}
        \item Evaluasi Verifikator: Kalimat merangkai sempurna (Tertutup). Secara historis bernilai \textbf{SALAH (F)}.
    \end{itemize}

    \item \textbf{Subjek Konstanta Kedua:} $b = \text{Megawati Soekarnoputri}$
    \begin{itemize}
        \item Substitusi: $P(b)$
        \item Formasi Utuh: \textit{"Megawati Soekarnoputri adalah Presiden Pertama Republik Indonesia."}
        \item Evaluasi Fakta Dokumen: Menghasilkan vonis \textbf{SALAH (F)}.
    \end{itemize}
    
    \item \textbf{Subjek Konstanta Pamungkas:} $s = \text{Ir. Soekarno}$
    \begin{itemize}
        \item Substitusi Akhir: $P(s)$
        \item Kalimat berbunyi merdu: \textit{"Ir. Soekarno adalah Presiden Pertama Republik Indonesia."}
        \item Evaluasi Fakta Final: Memperoleh mahkota kebenaran sakral bernilai \textbf{BENAR (T)}.
    \end{itemize}
\end{enumerate}
\end{contoh}

Melalui runtutan eksperimen di atas, kita memaknai kaidah tertinggi FOL pertama kali: Sebagaimana mesin cetak gol \textit{striker} sepakbola takkan mengkonversi poin tanpa umpan cantik bola, maka **Predikat takkan pernah sanggup melontarkan kesimpulan Benar/Salah sebelum ada Subjek Variabel yang disubstitusi tuntas membumi di tangannya**. Transformasi \textbf{Kalimat Terbuka $\longrightarrow$ Proposisi} mensyaratkan penuangan Konstanta definitif padanya.
